\subsection{p-群}

在本节及下一节中，字母 \(p\) 表示一个素数（§ 4，第10号，命题16）。

定义 9. 阶为 p 的幂的有限群称为 p-群。

设 G 是一个阶为 \(p^{r}\) 的 p-群。\(p^{r}\) 的每个因子都是 p 的幂（ \(\S 4\) ，第 10 号，定理 7 的推论）。因此，G 的每个子群和每个商群都是 p-群（ \(\S 4\) ，第 4 号，命题 4 的推论）；G 的每个齐次空间的基数都是 p 的幂（ \(\S 5\) ，第 5 号，定理 1）。

p-群被 p-群的扩张是 p-群。

*例子。(1) 交换 \(p\) -群同构于循环群 \(\mathbf{Z} / p^n\mathbf{Z}\) 的乘积（参见习题 19，以及 VII, §4 第 7 号命题 7）。

\begin{enumerate}
\def\labelenumi{(\arabic{enumi})}
\setcounter{enumi}{1}
\item
  设 \(k\) 为特征 \(p\) 的有限域。严格三角群 \(T_1(n, k)\) 是一个 \(p\) -群。
\item
  四元数群 \(\{\pm 1, \pm i, \pm j, \pm k\}\) 是一个2-群（参见习题4）。*
\end{enumerate}

命题 11. 设 E 为有限集，G 为作用在 E 上的 p-群。用 \(E^{G}\) 表示由 \(x \in E\) （满足对所有 \(g \in G\) 有 gx = x）所成的集合（不动点）。则

\[
\operatorname{Card} \left(\mathrm{E} ^ {\mathrm{G}}\right) \equiv \operatorname{Card} (\mathrm{E}) \quad (\text { mod. } p).
\]

\(E - E^{G}\) 是一些不缩减为一点的轨道的不交并。这种轨道的基数是异于 \(p^{0} = 1\) 的 p 的幂，因而是 p 的倍数。

推论. 设 G 是一个 p-群。若 G 不缩减为 e，则其中心也不缩减为 e。

设 G 经内自同构作用在自身上。不动点集是 G 的中心 Z。由命题 11，

\[
\operatorname{Card} (Z) \equiv \operatorname{Card} (G) \equiv 0 \pmod {p},
\]

因此 \(\operatorname{Card}(\mathbf{Z}) \neq 1\) ，且 \(\mathbf{Z} \neq \{e\}\) 。

定理 1. 设 G 是一个 p-群，其阶为 \(p^{r}\) 。存在 G 的子群序列

\[
\mathbf {G} = \mathbf {G} ^ {1} \supset \mathbf {G} ^ {2} \supset \dots \supset \mathbf {G} ^ {r + 1} = \{e \}
\]

使得 \((\mathbf{G},\mathbf{G}^k)\subset \mathbf{G}^{k + 1}\) （ \(1\leqslant k\leqslant r\) ），且 \(\mathbf{G}^k /\mathbf{G}^{k + 1}\) （ \(1\leqslant k\leqslant r\) ）是阶为 \(p\) 的循环群。

该定理对 \(G = \{e\}\) 成立。我们对 \(\operatorname{Card}(G)\) 作归纳来证明它。设 Z 为 G 的中心， \(x \neq e\) 是 Z 中的一个元素（命题 11 的推论）， \(p^{s}, s \neq 0\) 是 x 的阶。则 \(x^{p^{s-1}}\) 是阶为 p 的元素，因此 Z 包含一个子群 \(G^{r}\) ，它是 p 阶循环群。根据归纳假设，群 \(G' = G/G^{r}\) 有一个满足所需性质的子群列 \((\mathbf{G}'^{k})_{1 \leqslant k \leqslant r}\) 。设 \(\pi: G \to G'\) 为典范同态。由 \(G^{k} = \pi^{-1}(\mathbf{G}'^{k})\) （ \(1 \leqslant k \leqslant r\) ）， \(G^{r+1} = \{e\}\) 所定义的 G 的子群序列是一个解，因为 \(G^{k}/G^{k+1}\) 同构于 \(G'^{k}/G'^{k+1}\) （ \(1 \leqslant k \leqslant r\) ）（§4，第 7 号，定理 4）。

推论. 每个 p-群都是幂零群。

这由第 3 号命题 7 得出。

命题 12. 设 G 是一个 p-群，H 是 G 的一个不同于 G 的子群。则：

\begin{enumerate}
\def\labelenumi{(\alph{enumi})}
\item
  正规化子 \(\mathbf{N}_{\mathbf{G}}(\mathbf{H})\) （ \(\mathbf{H}\) 在 \(\mathbf{G}\) 中的）不同于 \(\mathbf{G}\) 。
\item
  存在 G 的一个正规子群 N，其在 G 中的指数为 p，且包含 H。
\end{enumerate}

断言 (a) 由第 3 号命题 8 的推论 1 得出。我们证明 (b)。由第 3 号命题 8 的推论 2，存在正规子群 \(\mathbf{N}'\) （ \(\mathbf{G}\) 的），包含 \(\mathbf{H}\) ，不同于 \(\mathbf{G}\) ，且使得 \(\mathbf{G} / \mathbf{N}'\) 是交换的。设 \(\mathbf{N}\) 是一个不同于 \(\mathbf{G}\) 且包含 \(\mathbf{N}'\) 的极大子群。则 \(\mathbf{N}\) 是正规的（第 2 号，命题 6 的推论 3），且 \(\mathbf{G} / \mathbf{N}\) 是单交换 \(p\) -群，从而是阶为 \(p\) 的循环群（ \(\S 4\) ，第 10 号，命题 20 的推论）。

推论. 设 G 是一个 p-群。G 的每个指数为 p 的子群都是正规子群。

